\section{7. Autonomous services and data pipeline}

\subsection{7.1 The autonomous service}

The controller is rule-based feedback control with a learned feed-forward
term. It was chosen as the simplest thing that meets the requirements: the
rules can be checked by reading them, a decision is a handful of arithmetic
operations against a budget of one real second per model minute, and the
learned part is small enough to inspect by eye.

\textbf{Inputs.} Per room: the latest \COtwo{}, temperature and people count,
each with its model time; the room's sizing and its actuators' addresses, from
the registry; and, for the predictive policy, the occupancy profile from
storage. The controller never sees the true state of a room.

\textbf{Ventilation.} Each \COtwo{} reading produces one airflow command. Two
terms are combined:
\begin{align*}
q_{\mathrm{reactive}}(C) &= q_{\mathrm{empty}} + \mathrm{clamp}\!\left(\frac{C - 600}{1000 - 600},\,0,\,1\right)\,(q_{\max} - q_{\mathrm{empty}}), \\
q_{\mathrm{people}}(n) &= \min\bigl(7\,n + 0.35\,A,\; q_{\max}\bigr)\ \text{l/s}\quad\text{for } n \geq 0.5,\ \text{otherwise } 0.
\end{align*}
The reactive term rises from the unoccupied minimum at 600~ppm to the maximum
at 1000~ppm. The people term is the workplace rule for $n$ people in a room of
area $A$. The four policies differ only in how they use them:

\begin{longtable}[]{@{}L{0.16\linewidth}L{0.76\linewidth}@{}}
\toprule
\textbf{Policy} & \textbf{Airflow asked for} \\
\midrule
\endhead
\code{constant} & Design flow, always \\
\code{reactive} & $q_{\mathrm{reactive}}(C)$ \\
\code{predictive} & $\max\bigl(q_{\mathrm{reactive}}(C),\ q_{\mathrm{people}}(n)\bigr)$, where $n$ is the larger of the people counted now and the people the profile expects within the lead time \\
\code{oracle} & The same, with $n$ the true peak occupancy within the lead time, replayed from a recorded run of the same dates \\
\bottomrule
\end{longtable}

The lead time is the room's own \COtwo{} time constant at design flow, kept
between 10 and 30 minutes: 12 minutes for a lecture room, 22 for the fika room
1570, 30 for an office. The controller asks for as little as the policy wants
and leaves the floors and the override to the damper.

\textbf{Heating} is the same under every policy, so that the policies differ
only in ventilation: 21~°C on weekdays 07--18 and an 18~°C setback otherwise,
with an optimal start. The setback ends early enough for the room to be warm at
07:00: the degrees still missing, divided by an assumed warm-up rate of
1.5~K/h, plus 15 minutes for the valve and the sensor. A room that has started
pre-heating keeps going until 07:00 even if it warms faster than assumed. It is
the same idea as the predictive ventilation, applied to the slower process: the
schedule is known, so act ahead of it.

\textbf{The learned part} is a per-room occupancy profile by time of day
(\code{internal/predict}), trained offline by the \code{train} job on the
people counters' stored readings, never on the truth.
\begin{itemize}
\tightlist
\item A weekday is cut into 96 slots of 15 minutes.
\item A day's value for a slot is the 75th percentile of that slot's readings,
  so a count has to show for about four of the fifteen minutes to register.
  A gathering does; the counter's stray ``one person'' does not.
\item The profile keeps, per slot, the 80~\% quantile of those values across
  the training days. A day with no reading in a slot counts as empty.
\item The forecast for ``the next lead time'' is the highest of the slots the
  window touches. Below half a person it counts as an empty room.
\end{itemize}
The controller ventilates for a high quantile on purpose: a fan running a
little early costs less than catching up on stale air. The profile was kept
this simple because the evaluation shows the model is not what limits it
(Section 11): it foresees the fika breaks, and no model could foresee lectures
that land in a random room each day without the timetable as an input.

\textbf{When the data is thin.} With no trained model the predictive policy
acts as the reactive one and says so (\code{no\_forecast}). The controller
looks for a newer model every 5 minutes real and keeps the last one it loaded
if storage is unreachable.

\textbf{When an input cannot be trusted.} Every decision record names the mode
it was made in:

\begin{longtable}[]{@{}L{0.22\linewidth}L{0.36\linewidth}L{0.34\linewidth}@{}}
\toprule
\textbf{Mode} & \textbf{When} & \textbf{What is commanded} \\
\midrule
\endhead
\code{normal} & Every input is trusted & The policy's airflow; the heating schedule \\
\code{no\_forecast} & No trained model yet & \code{predictive} acts as \code{reactive} \\
\code{degraded\_occupancy} & The people count is older than 180~s, missing, or found stuck (\code{predictive} only; no other policy uses the count) & The airflow for the forecast and \COtwo{}, without the count \\
\code{degraded\_co2} & The \COtwo{} sensor is found stuck & Design flow \\
\code{degraded\_temp} & The temperature sensor is found stuck & 21~°C \\
\bottomrule
\end{longtable}

A \emph{silent} \COtwo{} or temperature sensor needs no mode. Decisions are
made per reading, so its room gets no commands, the actuator's last command
expires after 300~s, and the actuator goes to design flow or 21~°C by itself.
A \emph{stuck} sensor keeps reporting on time, so the controller has to notice
(D-11):
\begin{itemize}
\tightlist
\item \textbf{\COtwo{} and temperature} readings are noisy and rounded, so a
  healthy sensor seldom repeats itself. The limit is the shortest run of
  identical readings that a healthy 50-room building produces less than once
  in 100 days: 5 readings for \COtwo{}, 12 for temperature. One different
  reading and the sensor is trusted again.
\item \textbf{The people counter} cannot be judged that way, because an empty
  room really does read 0 for hours. It is checked against the room's \COtwo{}
  with two bounds from the mass balance of Section 6 that hold for any
  airflow. \COtwo{} cannot rise faster than the counted people can make it
  rise; if it does, more people are there. And after $n$ people have been
  there for a time $t$, \COtwo{} cannot be lower than they would give with the
  damper fully open all that time; if it is, fewer are there. Both are tested
  on five-minute means, with a margin of five standard errors, and allow a
  person twice or half the nominal \COtwo{}.
\end{itemize}

\textbf{When the controller itself is wrong, slow or gone.} The damper enforces
the limits on its own (D-6). It refuses a command outside the room's range. It
never goes below the flow floor. It opens fully while the latest \COtwo{}
reading is above 1100~ppm, taking that reading from the broker itself, not from
the controller. And when no valid command has arrived for 300~s it returns to
design flow. A dead controller therefore costs energy, not air quality.

\textbf{Outputs.} One command per reading to the room's actuator (IF-8) and
one decision record (IF-6) with the inputs, the mode, the target and what the
actuator answered.

\subsection{7.2 The data pipeline}

\textbf{Ingestion.} A sensor device writes each reading to BuildSim and then
publishes it on MQTT with its model time and a sequence number. The broker is
a relay, not a buffer: it keeps nothing on disk and its clients do not use
persistent sessions. A consumer that was away has missed what was sent, and
the gap stays visible. We chose that over buffering because a ventilation
decision made on a reading several minutes old is worse than no decision,
which the actuators already handle.

\textbf{Storage.} The storage service subscribes to all four record streams
and appends each record to a JSONL file per run and stream
(\code{<run>/obs.jsonl}, \code{truth.jsonl}, \code{decision.jsonl},
\code{event.jsonl}). Nothing is aggregated and nothing is deleted: a run of
five simulated weekdays is about 2.2 million records. One process owns the
files and serves them over HTTP (IF-10), because processes may share no files.
A time-series database was considered and rejected for this size (D-4). What
the control loop needs from the recent past, the last few minutes per room, it
keeps in its own memory.

\textbf{Processing.} Three consumers turn stored records into something else,
all off the control path. \code{train} reads a run's people counts and writes
the profile. \code{cmd/evaluate} reads the truth of finished runs and produces
every table in Section 11. viz keeps one point per room and model minute for a
week in memory, about 16~MB for 50 rooms, because a page that redraws every two
seconds cannot wait for storage to read a whole run; after a restart it reads
the run so far from storage once.

\textbf{Data quality.} A good decision on bad data is still a bad outcome, so
each kind of bad data has a defined handling:

\begin{longtable}[]{@{}L{0.20\linewidth}L{0.72\linewidth}@{}}
\toprule
\textbf{Problem} & \textbf{Handling} \\
\midrule
\endhead
Duplicate & QoS~1 can deliver twice. The controller and storage drop a repeated \code{(sensor\_id, seq)} \\
Stale & Judged by the reading's model time, never by its arrival: older than 180~s is not acted on \\
Missing & No reading, no decision; the actuator's command expires and it falls back. The dashboard raises ``no reading'' \\
Stuck & Found by the detector above and then treated like missing. The decision names the fault \\
Malformed & Storage rejects the record, counts it and logs it. The count is visible at \code{/healthz} \\
Gaps in the truth & Not filled. The evaluation counts the minutes it has \\
Wrong run & Every record carries the run id of the process that wrote it, and storage files it by that. A single service restarted under another run id would file its records there, so the deploy scripts take the run id from the running controller \\
\bottomrule
\end{longtable}

\textbf{Provenance.} Every number in Section 11 is produced by
\code{cmd/evaluate} from the stored records of a named run, and a model
records which runs and how many readings it was trained on.
